\subsection{A localization gate with one, two or three disks per weight}

Let $P_r(N)$ denote the specified first skeleton with $r$ copies
of each disk in (SK.1), for every weight $w\geq N$, and with
no generators below $N$. All copies have the same endpoints
$0,T^{[w]}$. Here $r$ is constant across the weights of one
object. For fixed $d\geq0$, let $\mathcal F_3^{[d]}$ be the
full subcategory of $\mathcal C^{[d]}$ whose first skeleton is
strictly isomorphic to some $P_r(N)$ with $N\geq2$ and
$r\in\{1,2,3\}$, and whose rows in form degrees through
$d+2$ are exact in every weight at least $N$.
The global coefficient comparison through degree $d$ is
already part of $\mathcal C^{[d]}$.
Let $\mathcal F_{3,\omega}^{[d]}$ be its full subcategory of
countable objects, in the precise sense of A-S09-29.

\begin{theorem}[A-S09-30: Proved; a three-disk terminal-localization gate]
\label{AS0930}
The actual inclusions
\[
 \mathcal F_{3,\omega}^{[d]}
       \longrightarrow\mathcal F_3^{[d]}
       \longrightarrow\mathcal A_{\mathbb Z}
 \tag{S09.161}
\]
induce equivalences on all-arrow localizations.
Consequently the terminal-localization assertion for
$\mathcal A_{\mathbb Z}$ is equivalent to terminality for
$\mathcal F_{3,\omega}^{[d]}$.
The inclusions
\[
 \mathcal E_\omega^{[d]}
       \subset\mathcal F_{3,\omega}^{[d]}
       \subset\mathcal C_\omega^{[d]}
 \tag{S09.162}
\]
are both proper, even up to strict isomorphism.
All statements hold with support $q\leq w$, termwise
integral flatness, or both.

The common-source gate remains the smaller
$\mathcal E_\omega^{[d]}$ of A-S09-29. This theorem proves
no terminality assertion for the new gate and makes no
claim that two copies alone suffice for this localization
comparison.
\end{theorem}

\begin{proof}
We give every comparison used in the localization.
For a finite collection of admissible objects $A$, choose
one integer
\[
 N\geq\max_A\{2,\lambda_A(d+2),\mu_A(d+2)\}.
 \tag{S09.163}
\]
For each such object, $T_NA$ has the coefficient comparison
through $d+2$ in every weight and has exact horizontal rows
through $d+2$ in weights at least $N$.
A chosen coefficient-edge system is an actual strict map
$p:P_1(N)\to(T_NA)_{\leq1}$, obtained by choosing one
normalized edge with zeroth face $T^{[w]}$ for each $w\geq N$.
Row-zero exactness supplies these elements.

For an ordered list of $r\leq3$ such systems, define
\[
 R_A(p_1,\ldots,p_r)
   =T_NA\mathop{\times}_{C((T_NA)_{\leq1})}
                  C(P_r(N)).
 \tag{S09.164}
\]
The evaluation of the $i$th copy in the right factor is $p_i$.
Write $\pi_A(p_1,\ldots,p_r)$ for its projection to $T_NA$
followed by the actual tail map to $A$.
Exactly the graph and pullback sequences of A-004 apply:
the free disk algebra has its integral coefficient contraction,
and its endpoint map is onto from weight $N$.
They give horizontal exactness through $d+2$ from $N$
and vertical coefficient comparison through $d$ above $N$.
Below $N$ every term of this pullback is literally $S$.
Thus (S09.164) belongs to $\mathcal F_3^{[d]}$.
These are the same uniform estimates independently checked
in A-S09-27, with $r$ disk copies in place of one.
All maps preserve support, and the finite matching subgroups
and pullbacks remain torsionfree when the original terms are
flat. No bound depends on the simplex degree.

Several elementary maps between (S09.164) will be used.
Inclusion of chosen copies of $P_r(N)$ into $P_s(N)$ gives
the corresponding strict map of replacements, with the
identity on $T_NA$. For an actual map $f:A\to B$, let
$fp$ abbreviate $(T_Nf)_{\leq1}p$.
The map $T_Nf$, together with the identity on the ordered
free generators, gives an actual map
\[
 R_f(p_1,\ldots,p_r):R_A(p_1,\ldots,p_r)
                         \longrightarrow R_B(fp_1,\ldots,fp_r).
\]
The equality of their evaluations verifies that the pair
lies in the target pullback. These assignments preserve
identity and composition as literal maps. They commute
strictly with inclusion of copies and with the projections
$\pi$. A fold of repeated copies is likewise an actual map
whenever their evaluations agree.

For two edge systems in the same object define the following
morphism in the localization of $\mathcal F_3^{[d]}$:
\[
 c_A(p,q)=[i_q]^{-1}[i_p]:R_A(p)\longrightarrow R_A(q),
 \qquad
 R_A(p)\xrightarrow{i_p}R_A(p,q)\xleftarrow{i_q}R_A(q).
 \tag{S09.165}
\]
Both arrows in this definition are the actual copy inclusions.
We prove the exact identities needed below.

First, folding $P_2(N)$ to $P_1(N)$ when both evaluations
are $p$ gives a strict map $r:R_A(p,p)\to R_A(p)$ with
$r i_1=r i_2=\mathrm{id}$. Inverting $r$ gives
$c_A(p,p)=\mathrm{id}$.
Second, for three systems $p,q,t$, use the actual object
$R_A(p,q,t)$. Each two-copy object maps to it by inclusion
of the indicated copies. Write $j_p,j_q,j_t$ for its
three maps from the corresponding one-copy objects.
The commutative inclusion triangles give
\[
 c_A(p,q)=[j_q]^{-1}[j_p],\qquad
 c_A(q,t)=[j_t]^{-1}[j_q],\qquad
 c_A(p,t)=[j_t]^{-1}[j_p].
\]
Cancellation therefore proves
$c_A(q,t)c_A(p,q)=c_A(p,t)$.
Finally the actual map $R_f(p,q)$ and its two inclusion
squares prove
\[
 c_B(fp,fq)[R_f(p)]
                       =[R_f(q)]c_A(p,q).
\]
All objects used in these proofs have at most three copies.
These identities also imply inversion of $c_A(p,q)$ by
$c_A(q,p)$. Every equality has been proved by the displayed
strict maps before formal inverses are taken.

For the finite collection of objects choose one edge system
$p_A$ in each. On the full category with those objects define
an actual functor to the localized category by
\[
 \begin{split}
 \Phi_N(A)&=R_A(p_A),\\
 \Phi_N(f)&=c_B(fp_A,p_B)[R_f(p_A)].
 \end{split}
 \tag{S09.166}
\]
The target is a groupoid; the maps in (S09.166) are specified
words in actual arrows of $\mathcal F_3^{[d]}$.
The identity law follows from the fold identity above.
For $A\xrightarrow f B\xrightarrow g C$, the naturality
identity for $c$ and then its three-copy identity give
\[
 \begin{split}
 \Phi_N(g)\Phi_N(f)
 &=c_C(gp_B,p_C)c_C(gfp_A,gp_B)
                         [R_g(fp_A)][R_f(p_A)]\\
 &=c_C(gfp_A,p_C)[R_{gf}(p_A)]
 =\Phi_N(gf).
 \end{split}
\]
This proves functoriality even if the full category on the
chosen finite object set has infinitely many morphisms.
No functorial choice of a numerical cutoff for the entire
admissible category has been asserted.

The projections compare this functor with the original arrows.
For every pair of systems $p,q$ in $B$, the two projections
from the copy-inclusion diagram satisfy the strict equalities
$\pi_B(p,q)i_p=\pi_B(p)$ and
$\pi_B(p,q)i_q=\pi_B(q)$.
Consequently (S09.165) and the projection square for $R_f$
give, in the ambient localization,
\[
 [\pi_B(p_B)]\,\Phi_N(f)=[f]\,[\pi_A(p_A)].
 \tag{S09.167}
\]
When both $A$ and $B$ belong to $\mathcal F_3^{[d]}$,
this equality already holds in its own localization:
the original arrow, the projections to $A,B$, and every
one-copy or two-copy intermediate are actual arrows between
objects of that full subcategory. This is the specific
internal equality that a single-copy bridge did not supply
in A-S09-28; here it follows from the strict same-base
projection diagram.

We now prove the localization equivalence without an infinite
cutoff choice. Every admissible object receives a projection
$R_A(p_A)\to A$ from an object of $\mathcal F_3^{[d]}$
for a sufficiently large $N$ in (S09.163), proving essential
surjectivity. For a word of ambient arrows and their inverses
between objects $X,Y\in\mathcal F_3^{[d]}$, choose one $N$
for its finitely many objects. Applying (S09.166) to its
arrows gives an internal word $\Phi_N(w)$, and
\[
 [\pi_Y(p_Y)]\,\Phi_N(w)\,[\pi_X(p_X)]^{-1}
\]
is an internal word whose ambient image is $w$, by (S09.167)
on each arrow and its inverse. This proves fullness.

For faithfulness, take two internal words $u,v$ that become
equal in the ambient all-arrow localization. An equality in
its presentation has a finite witness made from composition,
identity, and formal inverse cancellation relations.
Choose one $N$ for all objects in that finite witness and one
system $p_A$ for each of them. The functor (S09.166) sends
every witness relation to a valid equality in the localized
$\mathcal F_3^{[d]}$, so $\Phi_N(u)=\Phi_N(v)$ there.
On the objects and arrows of the original internal words,
(S09.167) is an internal equality, as just proved.
Conjugating by their endpoint projections gives $u=v$
in that localization. This proves faithfulness and the
second equivalence in (S09.161), in every optional variant.

For the first equivalence use the full finite-diagram
construction of A-S09-29, with the following exact
specialization. At each vertex of a finite diagram in
$\mathcal F_3^{[d]}$, include its whole first skeleton
$P_r(N)$ at the start; it is countable. Retain its actual
row cutoff $N$. Close countable sets under every indicated
algebra and simplicial operation and every specified arrow,
and adjoin the same vertical and normalized horizontal
primitives in (S09.157)--(S09.158). Their availability and
all resulting exactness ranges are exactly those proved
there. The first skeleton stays equal to $P_r(N)$, since
it was already included and all subsequent elements are
chosen inside the original object. The resulting countable
subdiagram therefore lies in $\mathcal F_{3,\omega}^{[d]}$.
Specified countable vertices stay fixed. The finite-word
and finite-equality-witness proof in A-S09-29 applies to
these very inclusions and proves the first equivalence.
Subgroup preservation of support and integral flatness
is unchanged. Composing the two equivalences proves the
terminal-localization reduction in the theorem.

We finish with strict properness. The doubled-skeleton object
$D$ constructed in A-S09-28 belongs to
$\mathcal F_{3,\omega}^{[d]}$. Its first nonzero
simplex-one zero-form quotient by coefficients is in
weight $N$ and has rank two. In every $P_1(M)$ the
first such quotient has weight $M$ and rank one.
A coefficient-preserving graded isomorphism would force
$M=N$ and then identify these distinct free ranks.
Thus $D$ is outside $\mathcal E_\omega^{[d]}$.

For the other inclusion take $C_*=R(T_NG)$ with $N$
as in (S09.144). It lies in $\mathcal C^{[d]}$ by
A-S09-24 and is countable: each standard simplex group
is countable, each filler set is a subset of such a
group, the union over its countably many weights is
countable, and every coskeleton matching diagram and
pullback in a fixed simplex is a finite limit of
countable groups. Its actual first weight-$N$ filler
family contains every distinct edge (S09.145).
The first skeleton supplied by A-006 therefore has
infinitely many independent disk variables in weight
$N$ and none below $N$. Its first nonzero quotient
by coefficients has infinite free rank. Every
$P_r(M)$ allowed here has first nonzero quotient in
weight $M$ of rank $r\leq3$. Such an isomorphism
would force $M=N$ and equate infinite rank with one
of these three finite ranks, which is impossible
already after tensoring with $\mathbb Q$.
Hence $C_*\notin\mathcal F_{3,\omega}^{[d]}$.
Both properness witnesses are flat and satisfy the
support condition, as checked in their constructions.
This proves (S09.162) and all remaining assertions.
\end{proof}
